\section{Main Theorem for Type $A$ Springer Resolution}
\label{chapter_main_thm}

Consider the self-dual of symplectic resolution $X = X^! = \tilde{\nilcone} \to X_0 = X_0^! = \nilcone$ given by the Springer resolution of type $A$. We would like to apply the formalism in \Cref{sec:bichar_poly} and compute the bicharacteristic polynomial.

We will prove that the bifiltration on cohomology is the direct sum of bifiltrations on the summands of \Cref{crl:type_A_decomposition}, which implies our main theorem \Cref{thm:bichar_full_flag}.

\begin{lemma}[lem:main_product_filtration]{Main Lemma}
    The bicharacteristic polynomial $T_{X,X^!}(x, y)$ equals the sum of
    \[
        x^{l(\lambda)}Q(\OO_\lambda, x^{-1})\cdot y^{l(\lambda^T)}Q(\OO_{\lambda^T}, y^{-1}) = K_{\lambda,(1)^n}(x)\cdot K_{\lambda^T,(1)^n}(y),
    \]
    over all partitions $\lambda$ of $n$.
\end{lemma}

\subsection{Category \mathintitle{\OO} and the BBD Summands}

We would like to show that the isomorphism $\misom : \HH_T(\tilde{\nilcone})_\sigma \to \HH_T(\tilde{\nilcone})_\eta$ in \Cref{eqn:m} preserves the BBD summands. To see this, we need to employ techniques from category $\mathcal{O}$. We will just briefly walk through the results we need.

By the general philosophy of \cite{braden2016quantizationsconicalsymplecticresolutions}, the pair of dominant integral characters $(\sigma, \eta)$ determines Koszul dual categories $\mathcal{O}$ and $\mathcal{O}^!$, which gives isomorphisms
\begin{equation*}
    K(\mathcal{O}) \xrightarrow{\sim} \HH_T(\tilde{\nilcone})_\sigma,
    \qquad
    K(\mathcal{O}^!) \xrightarrow{\sim} \HH_T(\tilde{\nilcone})_\eta.
\end{equation*}

The objects $P_w, L_w, S_w \in K(\mathcal{O})$ (resp. $P^!_{w^!}, L^!_{w^!}, S^!_{w^!} \in K(\mathcal{O}^!)$) are called the projective objects, simple objects, and the standard objects, indexed by the Weyl group $W$. It is known that each of these classes of objects forms a basis in the $K$-groups, and hence in the specialized cohomology. Moreover, the standard objects corresponds to the stable envelopes in our context. We have the following result compairing these objects and the BBD decomposition.

For each $w \in W$, write $\lambda_w$ the unique partition satisfying $G \cdot (\lie{n} \cap w \cdot \lie{n}) = \nilorbit_{\lambda_w}$. Recall that the ordering on the set of partitions satisfies $\nilorbit_\mu \subseteq \overline{\nilorbit_\lambda}$ if $\mu \leq \lambda$.

\begin{proposition}[prop:]{}
    The span of the BBD summands $\leq \lambda$ is generated by the classes $\{ [P_w] \mid \lambda_w \leq \lambda \}$. Dually, the span of the BBD summands $\geq \lambda$ is generated by the classes $\{ [L_w] \mid \lambda_w \geq \lambda \}$.
\end{proposition}
We now choose the local polarizations $\varepsilon$ and $\varepsilon^!$. Define $\varepsilon_w = (-1)^{\ell(w)}$ and $\varepsilon^!_{w^!} = 1$ for all $w \in W$, where $\ell$ is the length function on $W$ (for type $A$, $(-1)^{\ell(w)}$ is just the sign of the permutation $w$). This choice is motivated by \cite[Remark 10.28]{braden2016quantizationsconicalsymplecticresolutions} as follows. Note that we have 
\begin{equation*}
    \misom([S_w]) = (-1)^{\ell(w)} \misom(\Stab_\varepsilon(wB)) = (-1)^{\ell(w)} \Stab_{\varepsilon^!}(w^!B) = (-1)^{\ell(w)} [S^!_{w^!}].
\end{equation*}
By the argument in \cite[Remark 10.28]{braden2016quantizationsconicalsymplecticresolutions}, we have that $\misom([L_w]) = (-1)^{\ell(w)} [P^!_{w^!}]$ and $\misom([P_w]) = (-1)^{\ell(w)} [L^!_{w^!}]$. It follows that the BBD summands $\leq \lambda$ in $\HH_T(\tilde{\nilcone})_\sigma$ is sent to the BBD summands $\geq \lambda^T$ in $\HH_T(\tilde{\nilcone})_\eta$ under the isomorphism $\misom$, and vice versa. By taking the associated graded basis, we arrived to the following result.

\begin{proposition}[prop:m_preserves_summand]{}
    The isomorphism $\misom : \HH_T(\tilde{\nilcone})_\sigma \to \HH_T(\tilde{\nilcone})_\eta$ sends the BBD summand corresponding to $\lambda$ to the BBD summand corresponding to $\lambda^T$.
\end{proposition}
    


% \begin{definition}[def:]{}
%     Let $\mathbb{F}$ be the fraction field of $\HH^\bullet_T(\pt)$. Define the intersection pairing
%     \[
%         \AB{\ ,\ }_{\mathbb{F}}:\HH^\bullet_T(X,\mathbb{C})\otimes_{\HH^\bullet_T(\pt)}\HH^\bullet_T(X)\to \mathbb{F}
%     \]
%     as the following symmetric pairing:
%     \[
%         \AB{A,B}_{\mathbb{F}}=(-1)^{\dim_\mathbb{R} X/4}
%         \sum_{x\in X^T}^{}\frac{i^*_xA\cdot i^*_xB}{\ee(T_xX)}.
%     \]
% \end{definition}

% \begin{lemma}[lem:]{\cite[Section 6.2]{braden2016quantizationsconicalsymplecticresolutions}}
%     $\AB{[S_\alpha], [S_\beta]} = \AB{[L_\alpha], [P_\beta]}=\delta_{\alpha,\beta}$.
% \end{lemma}

% In \cite[Remark 10.28]{braden2016quantizationsconicalsymplecticresolutions}, the authors defiend a pairing $\AB{\ ,\ }_{*}$ between $K_0(\OO)_\mathbb{C}$ and $K_0(\OO^!)_\mathbb{C}$ characterized by any of the following:
% \begin{alignat*}{1}
%     \AB{[S_\alpha], [S^!_\beta]}_* &= \delta_{\alpha,\beta}\\
%     \AB{[L_\alpha], [L^!_\beta]}_* &= \delta_{\alpha,\beta}\\
%     \AB{[P_\alpha], [P^!_\beta]}_* &= \delta_{\alpha,\beta},
% \end{alignat*}

% \begin{lemma}[lem:]{}
%     We have $\AB{A,B}_*=\AB{A,\mm(B)}$.
% \end{lemma}
%\todo{The key fact is that the resolution of a simple module $L_\alpha$ by standard modules $S_\beta$ has almost the same shape as the composition series of the dual projective module $P_\alpha^!$ by standard modules. }

\subsection{A Linear Algebra Lemma}

We now prove a linear algebra statement needed for the proof of the main statement. We will work on the following abstract setting.

Let $A$ and $B$ be (not necessarily commutative) algebras over $\mathbb{C}$. In what follows, we write $\otimes$, $\mathrm{Hom}$ and $\mathrm{End}$ for $\otimes_\mathbb{C}$, $\mathrm{Hom}_\mathbb{C}$ and $\mathrm{End}_\mathbb{C}$, unless a subscript is provided.

\begin{lemma}[lem:tensor_hom]{}
    Let $M$ and $M'$ be $A$-modules and $N$ and $N'$ be $B$-modules, all of finite dimension over $\mathbb{C}$. Then, the natural map
    \begin{equation*}
        \mathrm{Hom}_A(M, M') \otimes \mathrm{Hom}_B(N, N') \to \mathrm{Hom}_{A \otimes B}(M \otimes N, M' \otimes N')
    \end{equation*}
    is an isomorphism.
    \begin{proof}
        We treat both sides as subspaces of the vector space $\mathrm{Hom}(M \otimes N, M' \otimes N')$. Since an $(A \otimes B)$-module homomorphism is nothing but a linear map that is both an $A$-module homomorphism and a $B$-modules homomorphism, we have
        \begin{equation*}
            \mathrm{Hom}_{A \otimes B}(M \otimes N, M' \otimes N') = \mathrm{Hom}_A(M \otimes N, M' \otimes N') \cap \mathrm{Hom}_B(M \otimes N, M' \otimes N').
        \end{equation*}
        By the finite-dimensionality of $N$ and $N'$, we have
        \begin{equation*}
            \mathrm{Hom}_A(M \otimes N, M' \otimes N') = \mathrm{Hom}_A(M, M') \otimes \mathrm{Hom}(N, N').
        \end{equation*}
        Similarly, by the finite-dimensionality of $M$ and $M'$, we have
        \begin{equation*}
            \mathrm{Hom}_B(M \otimes N, M' \otimes N') = \mathrm{Hom}(M, M') \otimes \mathrm{Hom}_B(N, N').
        \end{equation*}
        It follows that their intersection is $\mathrm{Hom}_A(M, M') \otimes \mathrm{Hom}_B(N, N')$.
    \end{proof}
\end{lemma}

\begin{lemma}[lem:commutes_implies_split]{}
    Let $V$ and $W$ be finite-dimensional vector spaces, $M$ be a simple $A$-module and $N$ be a simple $B$-module. Let $f : V \otimes M \to W \otimes N$ be an isomorphism of vector spaces so that the $B$-action on $V \otimes M$ induced by $f$ commutes with the original $A$-action. Assume that $\dim_\mathbb{C} M = \dim_\mathbb{C} W$ and $\dim_\mathbb{C} V = \dim_\mathbb{C} N$. Then, there are isomorphisms of vector spaces $f_1 : V \to N$ and $f_2 : M \to W$ such that $f = f_1 \otimes f_2$.
    \begin{proof}
        Since the induced $B$-action on $V \otimes M$ commutes with the original $A$-action, we have an algebra homomorphism $B \to \mathrm{End}_A(V \otimes M)$. By \Cref{lem:tensor_hom} and Schur's lemma,
        \begin{equation*}
            \mathrm{End}_A(V \otimes M) \cong \mathrm{End}(V) \otimes \mathrm{End}_A(M) = \mathrm{End}(V).
        \end{equation*}
        It follows that the $B$-action on $V \otimes M$ comes from a $B$-action on $V$. Similarly, the $A$-action on $W \otimes N$ comes from an $A$-action on $W$. By \Cref{lem:tensor_hom} again, we have
        \begin{equation*}
            f \in \mathrm{Hom}_{A \otimes B}(V \otimes M, W \otimes N) = \mathrm{Hom}_A(M, W) \otimes \mathrm{Hom}_B(V, N).
        \end{equation*}
        Note that every $A$-module homomorphism from the simple $A$-module $M$ to $W$ is either zero or injective. Since $\dim_\mathbb{C} M = \dim_\mathbb{C} W$ and $f$ is an isomorphism, we must have $M \cong W$ as $A$-modules and thus $\mathrm{Hom}_A(M, W) \cong \mathbb{C}$. Similarly, $V \cong N$ as $B$-modules and $\mathrm{Hom}_B(V, N) \cong \mathbb{C}$. Our claim follows.
    \end{proof}
\end{lemma}


\subsection{Commuting Actions}

% \begin{equation} \label{eqn:M}
%     \HH^{\bullet+2d}_T(X)_\loc \otimes_\mathbb{C} \HH^\bullet_{T^!}(\mathrm{pt})_\loc \xleftarrow[\sim]{\Stab_{\sigma, \varepsilon}^X} W^{\bullet, \bullet}_\loc \xrightarrow[\sim]{\Stab_{\eta, \varepsilon^!}^{X^!}} \HH^\bullet_T(\mathrm{pt})_\loc \otimes_\mathbb{C} \HH^{\bullet+2d^!}_{T^!}(X^!)_\loc
% \end{equation}
% Since $\sigma$ and $\eta$ are generic, we can further tensor with the evaluation $\mathbb{C}[\lie{t} \times \lie{t}^!]_\loc \to \mathbb{C}$ at the point $(\sigma, \eta) \in \lie{t}_\mathrm{reg} \times \lie{t}^!_\mathrm{reg}$ .
% \begin{equation} \label{eqn:m}
%     \HH_T(X)_\sigma \xleftarrow[\sim]{\Stab_{\sigma, \varepsilon}^X} W_{\sigma, \eta} \xrightarrow[\sim]{\Stab_{\eta, \varepsilon^!}^{X^!}} \HH_{T^!}(X^!)_\eta
% \end{equation}

The isomorphism $\misom : \HH_T(\tilde{\nilcone})_\sigma \to \HH_T(\tilde{\nilcone})_\eta$ in \Cref{eqn:m} sends the left $\HBMtop(\steinberg)$-action on $\HH_T(\tilde{\nilcone})_\sigma$ to a left $\HBMtop(\steinberg)$-action on $\HH_T(\tilde{\nilcone})_\eta$, giving us two actions on the latter space.

\begin{proposition}[prop:action_commute]{}
    The two actions by $\HBMtop(\steinberg)$ on $\HH_T(\tilde{\nilcone})_\eta$ commutes with each other.
    \begin{proof}
        We show this on the level of equivariant cohomology. By localization, we can embed both $\HH_T^\bullet(\tilde{\nilcone}) \otimes_\mathbb{C} \HH^\bullet_T(\mathrm{pt})$ and $\HH^\bullet_T(\mathrm{pt}) \otimes_\mathbb{C} \HH_T^\bullet(\tilde{\nilcone})$ into $\mathbb{C}[\lie{h} \times \lie{h}]^{\oplus W}$. Denote by $\{1_w\}_{w \in W}$ (or by $\{1^!_w\}_{w \in W}$) the standard basis corresponding to the fixed points. By \Cref{lem:W_action_on_stab}, the $W$-action induced from the former is given by $w * 1_{w'} = 1_{w'w^{-1}}$. Similarly, the $W$-action induced from the latter is given by
        \begin{equation*}
            w *^! 1_{w'} = w *^! 1^!_{{w'}^{-1} w_\circ} = 1^!_{{w'}^{-1} w_\circ w^{-1}} = 1_{w_\circ w w_\circ w'},
        \end{equation*}
        where we are manipulating the equations $w^! = w^{-1} w_\circ$ and $w = (w_\circ w^{-1})^!$. Thus, the $*$-action is given by right multiplication while the $*^!$-action is given by left multiplication, hence they commute.
    \end{proof}
\end{proposition}

\subsection{Proof of the Main Theorem}

\begin{proof}[Proof of \Cref{lem:main_product_filtration}]
    Consider the restriction of $\misom$ in \Cref{eqn:m} to the summand
    \begin{equation*}
        \misom : \mathrm{IH}^\bullet(\nilorbit_\lambda) \otimes \Htop(\borel_\lambda) \to \mathrm{IH}^\bullet(\nilorbit_{\lambda^T}) \otimes \Htop(\borel_{\lambda^T}),
    \end{equation*}
    as indicated in \Cref{prop:m_preserves_summand}. Note that $\dim \Htop(\borel_\lambda)$ is given by the number of irreducible components of $\borel_\lambda$, which by \Cref{thm:irred_springer} is given by the number of tableaux of shape $\lambda$. On the other hand, $\dim \mathrm{IH}^\bullet(\nilorbit_{\lambda^T}) = K_{\lambda^T, (1)^n}(1)$ by \Cref{thm:IH_kostka}, which by definition is the number of tableaux of shape $\lambda^T$. The two numbers are the same. On the other hand, the $\HBMtop(\steinberg) \cong \mathbb{C}[W]$ acts only on $\Htop(\borel_\lambda)$ and $\Htop(\borel_{\lambda^T})$, both irreducibly due to \Cref{thm:springer_correspondence}. Finally, by \Cref{prop:action_commute}, the induced $W$-action by $\misom$ commutes with the original one. Thus, we can apply \Cref{lem:commutes_implies_split} to conclude our claim.
\end{proof}